% Gedruckte S. 132--141 einschließlich unnummerierter Blätter.
\begingroup
\clearpage
\noindent132\hfill{\small\textsc{Nachlass.\quad Übertragung der Kugel auf die Ebene etc.}}\hfill\mbox{}\par\medskip
\begingroup\small
[im anderen Fall]
\[
\begin{aligned}
x'-x&=\varepsilon\sin\beta\\
y'+y&=\varepsilon\cos\beta,
\end{aligned}
\]
[dann ergibt sich, da $\tfrac12(e^y+e^{-y})=\cos iy=\sec Y$ ist:]
\[
\begin{aligned}
\sin\tfrac12D.e^{\frac12i(A+A')}&=\frac{\sin\{\frac12\delta(\cos\alpha+i\sin\alpha)\}}{\sqrt{(\sec Y\sec Y')}}\\[5pt]
\cos\tfrac12D.e^{\frac12i(A-A')}&=\frac{\cos\{\frac12i\varepsilon(\cos\beta+i\sin\beta)\}}{\sqrt{(\sec Y\sec Y')}}.
\end{aligned}
\]
[Nun gilt]
\[
\begin{aligned}
\log\sin u&=\log u-\tfrac16uu-\tfrac1{180}u^4-\tfrac1{2835}u^6\ldots\\
\log\cos u&=\phantom{\log u}-\tfrac12uu-\tfrac1{12}u^4-\tfrac1{45}u^6\ldots,
\end{aligned}
\]
[folglich ergibt sich für $u=\tfrac12\delta e^{i\alpha}$ beziehungsweise $\tfrac12i\varepsilon e^{i\beta}$:]
\[
\begin{aligned}
\tfrac12(A+A')&=\alpha-\tfrac1{24}\delta\delta\sin2\alpha-\tfrac1{2880}\delta^4\sin4\alpha-\tfrac1{181440}\delta^6\sin6\alpha\ldots\\
\tfrac12(A-A')&=\phantom{\alpha-}\tfrac18\varepsilon\varepsilon\sin2\beta-\tfrac1{192}\varepsilon^4\sin4\beta+\tfrac1{2880}\varepsilon^6\sin6\beta\ldots,
\end{aligned}
\]
\[
\begin{aligned}
\log\sin\tfrac12D&=\log\tfrac12\delta-\tfrac1{24}\delta\delta\cos2\alpha-\tfrac1{2880}\delta^4\cos4\alpha-[\tfrac1{181440}\delta^6\cos6\alpha]\ldots\\
&\hspace{6em}-\tfrac12\log(\sec Y\sec Y')\\
\log\cos\tfrac12D&=\hspace{4em}\tfrac18\varepsilon\varepsilon\cos2\beta-\tfrac1{192}\varepsilon^4\cos4\beta+[\tfrac1{2880}\varepsilon^6\cos6\beta]\ldots\\
&\hspace{6em}-\tfrac12\log(\sec Y\sec Y').
\end{aligned}
\]
\begin{center}\rule{.18\textwidth}{.4pt}\par\medskip\textsc{Bemerkungen.}\end{center}
\begingroup\footnotesize
[1] ist, wie bereits auf S.~115 erwähnt, einem Notizbuch entnommen. Die Notiz [2] stammt aus einem anderen Notizbuch. Die Formeln für $R$ und $\dfrac{\partial R}{\partial y}$ sind dort berichtigt worden. Bei \textsc{Gauss} lauten sie:
\[
R=\tfrac12yy\operatorname{tang}\left(P-\frac xr\right)-\left(\tfrac14\operatorname{tang}\left(P-\frac xr\right)^2+\tfrac5{24}\right)y^4-\text{etc.}
\]
\[
\frac{\partial R}{\partial y}=\frac{\sin(P-x).(y+\tfrac16y^3+\tfrac1{120}y^5\ldots)}{\sqrt{(\cos(P-x)^2+2yy+\tfrac23y^4+\tfrac4{45}y^6\ldots)}}.
\]
Die Tafel [3] und die Notizen [4], [5], [6] und [8] sind verschiedenen Einzelblättern entnommen, auf denen sie zwischen Zahlenrechnungen aus den frühen 1820er-Jahren stehen. Die Notiz [7] wurde aus Aufzeichnungen in zwei Notizbüchern und auf einem einzelnen Zettel zusammengestellt. Ihre beiden Formeln für $\tfrac12(A+A')$ oben auf S.~130 wurden geändert. Im Original lauten sie:
\[
\begin{aligned}
\tfrac12(A+A')&=\alpha-\tfrac1{24}dd\sin2\alpha-\tfrac1{1440}d^4\sin4\alpha-\tfrac1{45360}.\tfrac1{64}d^6(2387\sin2\alpha+1039\sin6\alpha)\\
&=\alpha-\tfrac1{12}(x'-x)(y'-y)-\tfrac1{360}(x'-x)(y'-y)((x'-x)^2-(y'-y)^2).
\end{aligned}
\]
\hfill\textsc{Krüger, Börsch.}\par
\begin{center}\rule{.24\textwidth}{.4pt}\end{center}
\par\endgroup\par\endgroup
\clearpage
% Unnumbered opening in the position of printed 133.
\begin{center}
{\large STEREOGRAFISCHE DARSTELLUNG DES SPHÄROIDS\\{}[IN DER EBENE].}
\end{center}
\vspace{12mm}
\begin{center}
\begin{tabular}{@{}l@{ }l@{}}
$a\, .\, .\, .\, .\, .\, .\, .\, .\, .\, .\, .$&Große Halbachse\\[3pt]
$\varepsilon\, .\, .\, .\, .\, .\, .\, .\, .\, .\, .\, .$&Exzentrizität\\[3pt]
$90^\circ-W\, .\, .\, .\, .\, .$&Polhöhe des ersten Zentralorts\\[3pt]
$90^\circ-w\, .\, .\, .\, .\, .$&Polhöhe eines beliebigen Orts\\[3pt]
$\lambda\, .\, .\, .\, .\, .\, .\, .\, .\, .\, .\, .$&dessen westliche Länge\\[3pt]
$x+iy=t\, .\, .\, .$&[ebene] Koordinaten jenes Orts
\end{tabular}
\end{center}
\[
\left(\frac{1+\varepsilon\cos w}{1-\varepsilon\cos w}\right)^{\frac12\varepsilon}=\mathrm f(w),\qquad
\frac{2a}{\sqrt{(1-\varepsilon\varepsilon\cos w^2)}}=\mathrm F(w).
\]
\[
\begin{aligned}
\operatorname{tang}\tfrac12w&=b\\
\operatorname{tang}\tfrac12W&=B\\
\frac{e^{i\lambda}\mathrm f(w)}{\mathrm f(W)}&=h
\end{aligned}
\]
\[
\begin{array}{r@{\qquad}l}
1)&\displaystyle t=\frac{bh-B}{1+bhB}.\mathrm F(W).
\end{array}
\]
$x'+iy'=t'$ Koordinaten desselben Orts in Bezug auf den zweiten Zentralort;

$b^0$, $h^0$ haben dieselbe Bedeutung wie zuvor, jetzt in Bezug auf den zweiten Zentralort $[(90^\circ-w^0,\lambda^0)$, also $b^0=\operatorname{tang}\tfrac12w^0$, $h^0=\dfrac{e^{i\lambda^0}\mathrm f(w^0)}{\mathrm f(W)}]$.
\[
\begin{array}{r@{\quad}l}
2)&\displaystyle t'=\GaussSquareBracketPair{\mathrm F(w^0)\frac{\dfrac{bh}{h^0}-b^0}{1+\dfrac{bhb^0}{h^0}}=}
\mathrm F(w^0)\frac{1+b^0Bh^0}{b^0-Bh^0}.
\frac{t-\dfrac{b^0h^0-B}{1+b^0Bh^0}\mathrm F(W)}{t+\dfrac{h^0+b^0B}{b^0-Bh^0}\mathrm F(W)}.
\end{array}
\]
Vergrößerungsverhältnis:
\[
\sqrt{\frac{1-\varepsilon\varepsilon\cos w^2}{1-\varepsilon\varepsilon\cos W^2}}.
\frac{1+\dfrac{xx+yy}{(\mathrm F(W))^2}}{\cos\tfrac12w^2.\dfrac{\mathrm F(W)}{\mathrm f(w)}+\sin\tfrac12w^2.\dfrac{\mathrm F(w)}{\mathrm f(W)}}.
\]
\clearpage
\noindent134\hfill{\footnotesize\textsc{Nachlass.\quad Stereographische Darstellung des Sphäroids in der Ebene.}}\par\medskip
NB. Wenn man das Sphäroid schließlich stereografisch in der Ebene darstellen will, braucht man es nicht nach dem Verfahren auf S.~22 auf die Kugel zu übertragen. Dieses liefert
\[
x=A\operatorname{tang}\tfrac12(W-U)
\]
und lässt keine allgemeine Tafel für $U$ zu. Es genügt das Verfahren auf S.~21, bei dem $x$ die Form
\[
x=A'+B'\operatorname{tang}\tfrac12(W^*-U).
\]
erhält. Hier kann dann die allgemeine Tafel für $U$ verwendet werden.
\begin{center}\rule{.18\textwidth}{.4pt}\par\bigskip\textsc{Bemerkungen.}\end{center}
\begingroup\small
Die oben angeführten Seitenzahlen (21 und 22) beziehen sich auf die »Allgemeine Auflösung der Aufgabe: Die Theile einer gegebenen Fläche etc.«, die 1825 in Heft 3 der von H.~C.~\textsc{Schumacher} herausgegebenen »Astronomischen Abhandlungen« erschien. Gemeint ist Art.~13 dieser Abhandlung, Bd.~IV, S.~207 ff. Die obige Notiz selbst, die einem Einzelblatt entnommen ist, schließt an Art.~12, Bd.~IV, S.~206, an. Nach dieser Notiz wird die konforme Darstellung der Oberfläche des Sphäroids in der Ebene durch die Gleichung
\[
x+iy=\mathrm f\left(\lambda+i\log\left\{\operatorname{cotang}\tfrac12w\left(\frac{1-\varepsilon\cos w}{1+\varepsilon\cos w}\right)^{\frac12\varepsilon}\right\}\right)=\mathrm f(v)
\]
vermittelt. Für die stereografische Projektion ist aber
\[
\mathrm f(v)=k\operatorname{tang}(\operatorname{arc\,tang}e^{iv}-\tfrac12W)=k\frac{e^{iv}-\operatorname{tang}\tfrac12W}{1+e^{iv}\operatorname{tang}\tfrac12W}
\]
zu setzen, wobei $90^\circ-W$ die Breite des Punkts bezeichnet, der dem Augenpunkt diametral gegenüberliegt, also die Breite des Koordinatenursprungs. Für die Kugel gilt, wenn $2q$ und $2p$ die Komplemente der Breiten des Zentralorts und eines beliebigen Punkts sind, nach Notiz [3] auf S.~121:
\[
x+iy=k\frac{e^{iv}-\operatorname{tang}q}{1+e^{iv}\operatorname{tang}q},\qquad e^{iv}=e^{i\lambda}\operatorname{tang}p.
\]
\hfill\textsc{Krüger, Börsch.}\par
\begin{center}\rule{.24\textwidth}{.4pt}\end{center}
\par\endgroup
\clearpage
% Unnumbered title in the position of printed 135.
\null\vfill
\begin{center}
{\LARGE KONFORME ÜBERTRAGUNG\\[8pt]DES SPHÄROIDS\\[8pt]AUF DEN KEGELMANTEL.}
\end{center}
\vfill\null
\clearpage
% Blank source leaf in the position of printed 136.
\null
\clearpage
% Unnumbered opening in the position of printed 137.
\begin{center}
{\LARGE\textls[180]{NACHLASS.}}\par\vspace{13mm}
{\large [KONFORME ÜBERTRAGUNG DES SPHÄROIDS\\AUF DEN KEGELMANTEL.]}
\end{center}
Zur zweiten Darstellungsart des Sphäroids, auf einen Parallelkreis bezogen.
\begin{center}
\begin{tabular}{@{}l@{ }p{.72\textwidth}@{}}
$P\, .\, .\, .\, .$&Breite des Hauptparallelkreises\\[3pt]
$m\, .\, .\, .\, .$&Vergrößerungsverhältnis\\[3pt]
$\lambda\, .\, .\, .\, .\, .$&Länge eines beliebigen Punkts\\[3pt]
$\varphi\, .\, .\, .\, .$&dessen Breite\\[3pt]
$x,y\, .\, .$&Koordinaten seiner Darstellung in der Ebene [nach der Abwicklung des Kegelmantels auf die Ebene]\\
&\hspace{3em}$x=r\cos\Theta$, $y=r\sin\Theta$\\[3pt]
$R\, .\, .\, .\, .$&Wert von $r$ für $\varphi=P$
\end{tabular}
\end{center}
\[
\begin{aligned}
\mu&=\sin P\\
e\sin\varphi&=\sin u\\
e\sin P&=\sin U.
\end{aligned}
\]
\[
\begin{aligned}
\Theta&=\mu\lambda\\[4pt]
R&=\frac{a\cos P}{\mu\sqrt{(1-ee\sin P^2)}}=\frac{a\cos P}{\mu\cos U}=\frac{a}{\cos U\operatorname{tang}P}\\[5pt]
\frac rm&=\frac{a\cos\varphi}{\mu\sqrt{(1-ee\sin\varphi^2)}}=\frac{a\cos\varphi}{\mu\cos u}\\[5pt]
\frac{\mathrm dr}{r}&=-\frac{\mu(1-ee)\mathrm d\varphi}{(1-ee\sin\varphi^2)\cos\varphi}\\[5pt]
\frac{\mathrm dm}{m}&=\frac{\sin\varphi-\mu}{\cos\varphi}.\frac{1-ee}{1-ee\sin\varphi^2}\mathrm d\varphi=\frac{\mathrm dr}{r}\left(1-\frac{\sin\varphi}{\mu}\right)\\[5pt]
k&=\frac{a\operatorname{tang}(45^\circ+\frac12P)^\mu}{\cos U\operatorname{tang}(45^\circ+\frac12U)^{\mu e}\operatorname{tang}P}
\end{aligned}
\]
\vfill\noindent{\footnotesize IX.}\hfill18\hspace*{2cm}\par
\clearpage
\noindent138\hfill\textsc{Nachlass.}\hfill\mbox{}\par\medskip
\[
\begin{aligned}
r&=k\operatorname{tang}(45^\circ-\tfrac12\varphi)^\mu.\left(\frac{1+e\sin\varphi}{1-e\sin\varphi}\right)^{\frac12\mu e}\\
&=k\operatorname{tang}(45^\circ-\tfrac12\varphi)^\mu.\operatorname{tang}(45^\circ+\tfrac12u)^{\mu e}.
\end{aligned}
\]
Ist $\zeta$ das Azimut eines Elements einer geodätischen Linie gegenüber dem ersten Meridian, so ergibt sich
\[
\mathrm d\zeta=\left(1-\frac{\sin\varphi}{\mu}\right)\mathrm d\Theta=6q\,\mathrm d\Theta.
\]
Sind $q^0$, $q'$ die Werte von $q$ an den Endpunkten einer Geraden in der Ebene, deren Azimut $=Z$ ist, und $\zeta^0$ und $\zeta'$ die Azimute der geodätischen Linie an denselben Endpunkten, [ferner] $\Theta^0$, $\Theta'$ die [zu diesen gehörenden] Werte von $\Theta$, so ergibt sich
\[
Z=\zeta^0+(2q^0+q')(\Theta'-\Theta^0)=180^\circ+\zeta'+(q^0+2q')(\Theta^0-\Theta').
\]
Für $6q$ erhält man die folgenden Reihen, wenn man
\[
\log\operatorname{hyp}\frac rR=\rho\quad\text{und}\quad r=R(1+\delta)
\]
setzt:
\begingroup\footnotesize
\[
\begin{aligned}
{[}6q={]}&\frac{(1-ee\mu\mu)(1-\mu\mu)}{(1-ee)\mu\mu}\rho+\frac{(1+ee-2ee\mu\mu)(1-ee\mu\mu)(1-\mu\mu)}{(1-ee)^2\mu\mu}\rho\rho\\
&-\frac{(1+ee-3(1+4ee+e^4)\mu\mu+15ee(1+ee)\mu^4-14e^4\mu^6)(1-ee\mu\mu)(1-\mu\mu)}{3(1-ee)^3\mu^4}\rho^3\ldots\\[5pt]
=&\frac{(1-ee\mu\mu)(1-\mu\mu)}{(1-ee)\mu\mu}\delta+\frac{(1+3ee-4ee\mu\mu)(1-ee\mu\mu)(1-\mu\mu)}{2(1-ee)^2\mu\mu}\delta\delta\qquad\qquad .\\
&-\frac{(1+ee-(1+10ee+7e^4)\mu\mu+(9ee+21e^4)\mu^4-14e^4\mu^6)(1-ee\mu\mu)(1-\mu\mu)}{3(1-ee)^3\mu^4}\delta^3\ldots.
\end{aligned}
\]
\par\endgroup
Am bequemsten setzt man $q=\alpha(r-R)\sqrt r$ (wohl ebenso bequem $q=\alpha'r\log\frac rR$) und $\log m=\dfrac{\beta(r-R)^2}{\sqrt[3]r}$ und gibt die Logarithmen von $\alpha$ [und $\beta$] in einer Tafel an.

Für die Abplattung $=\dfrac1{302{,}78}$, d.~h. für das Achsenverhältnis $=\dfrac{301{,}78}{302{,}78}$, ergibt sich [vgl. S.~69]
\[
\begin{aligned}
\log ee&=7{,}819\,1850\,399\\
\log\sqrt{(1-ee)}&=9{,}998\,5632\,696\\
\log a&=6{,}804\,5975\,970.
\end{aligned}
\]
[Mit diesen Werten und für $P$ (Göttingen) $=51^\circ31'48''{,}7$ wurden die beiden folgenden Tabellen berechnet.]
\clearpage
\noindent\mbox{}\hfill{\small\textsc{Konforme Übertragung des Sphäroids auf den Kegelmantel.}}\hfill139\par\bigskip
\begin{center}
\begingroup\fontsize{6.7}{8.4}\selectfont\setlength{\tabcolsep}{2pt}\renewcommand{\arraystretch}{1.15}
\begin{tabular}{|r|r|r|r|r|r|r|r|}
\hline
$\varphi$&\shortstack{$\log m$\\(Ein-\\heiten d.\\7. Dezim.)}&$\log r$&\raisebox{0pt}[\dimexpr\height+2pt\relax]{\shortstack{$\log\dfrac1{\sqrt{(1-ee\sin\varphi^2)}}$\\(Einheiten d.\\7. Dezimalstelle)}}&\shortstack{$r$\\\\(Meter)}&$6q=1-\dfrac{\sin\varphi}{\sin P}$&$\log\alpha$&$\log\beta$\\\hline\hline
&&&&&&$-20$&$-20$\\[-3pt]
$47^\circ58'20''$&$8149.362$&$6{,}738\,2224.191$&$7915.864$&$547\,2961{,}838$&$0{,}051\,2376.087$&$8{,}964\,7419$&$7{,}962\,0287$\\
$48\ 25\ 0$&$6258.626$&$6{,}734\,2772.859$&$8026.658$&$542\,3470{,}546$&$0{,}044\,6331.717$&$8{,}964\,7802$&$7{,}962\,0627$\\
$48\ 51\ 40$&$4612.904$&$6{,}730\,2973.085$&$8137.252$&$537\,3993{,}625$&$0{,}038\,0862.194$&$8{,}964\,8096$&$7{,}962\,0902$\\
$49\ 18\ 20$&$3214.161$&$6{,}726\,2816.265$&$8247.620$&$532\,4534{,}268$&$0{,}031\,5971.467$&$8{,}964\,8304$&$7{,}962\,1111$\\
$49\ 45\ 0$&$2064.452$&$6{,}722\,2293.517$&$8357.734$&$527\,5083{,}660$&$0{,}025\,1663.435$&$8{,}964\,8430$&$7{,}962\,1257$\\
$50\ 11\ 40$&$1165.907$&$6{,}718\,1395.670$&$8467.569$&$522\,5640{,}957$&$0{,}018\,7941.974$&$8{,}964\,8476$&$7{,}962\,1336$\\
$50\ 38\ 20$&$520.757$&$6{,}714\,0113.250$&$8577.096$&$517\,6203{,}297$&$0{,}012\,4810.908$&$8{,}964\,8441$&$7{,}962\,1351$\\
$51\ 5\ 0$&$131.317$&$6{,}709\,8436\ 473$&$8686.291$&$512\,6767{,}793$&$0{,}006\,2274.051$&$8{,}964\,8333$&$7{,}962\,1302$\\
$51\ 31\ 40$&$0.004$&$6{,}705\,6355.223$&$8795.127$&$507\,7331{,}533$&$\underline{0{,}000\,0335.151}$&$8{,}964\,8152$&$7{,}962\,1188$\\
$51\ 58\ 20$&$129.337$&$6{,}701\,3859.037$&$8903.578$&$502\,7891{,}575$&$0{,}006\,1002.053$&$8{,}964\,7900$&$7{,}962\,1007$\\
$52\ 25\ 0$&$521.942$&$6{,}697\,0937.094$&$9011.617$&$497\,8444{,}952$&$0{,}012\,1733.878$&$8{,}964\,7581$&$7{,}962\,0761$\\
$52\ 51\ 40$&$1180\ 553$&$6{,}692\,7578.194$&$9119.218$&$492\,8988{,}664$&$0{,}018\,1856.670$&$8{,}964\,7202$&$7{,}962\,0460$\\
$53\ 18\ 20$&$2108.021$&$6{,}688\,3770.740$&$9226.356$&$487\,9519{,}681$&$0{,}024\,1366.806$&$8{,}964\,6761$&$7{,}962\,0107$\\
$53\ 45\ 0$&$3307.317$&$6{,}683\,9502.712$&$9333.004$&$483\,0034{,}926$&$0{,}030\,0260.713$&$8{,}964\,6263$&$7{,}961\,9676$\\
$54\ 11\ 40$&$4781.541$&$6{,}679\,4761.666$&$9439.136$&$478\,0531{,}307$&$0{,}035\,8534.839$&$8{,}964\,5713$&$7{,}961\,9186$\\
$54\ 38\ 20$&$6533.927$&$6{,}674\,9534.689$&$9544.727$&$473\,1005{,}674$&$0{,}041\,6155.684$&$8{,}964\,5115$&$7{,}961\,8636$\\
$55\ 5\ 0$&$8567.844$&$6{,}670\,3808.390$&$9649.754$&$468\,1454{,}848$&$0{,}047\,3209.773$&$8{,}964\,4470$&$7{,}961\,8019$\\[8pt]
$51\ 31\ 48{,}7$&$0$&$6{,}705\,6125.281$&$8795.718$&$507\,7062{,}715$&$0$&$8{,}964\,8150$&$7{,}962\,1184$\\
\end{tabular}\par\endgroup
\end{center}
\noindent\makebox[0pt][r]{.\hspace{2.5em}}[In der Spalte für $6q$ sind die Werte unterhalb des Strichs negativ.]
\vspace{12mm}
\begin{center}
\begingroup\fontsize{6.7}{8.4}\selectfont\setlength{\tabcolsep}{1.6pt}\renewcommand{\arraystretch}{1.12}
\begin{tabular}{|r|r|r||r|r|r||r|r|r||r|r|r|}
\hline
$\log r$&$\log\alpha$&$\log\beta$&$\log r$&$\log\alpha$&$\log\beta$&$\log r$&$\log\alpha$&$\log\beta$&$\log r$&$\log\alpha$&$\log\beta$\\\hline\hline
&$-20$&$-20$&&$-20$&$-20$&&$-20$&$-20$&&$-20$&$-20$\\[-3pt]
$6{,}670$&$8{,}964\,4414$&$7{,}961\,7965$&$6{,}687$&$8{,}964\,6612$&$7{,}961\,9980$&$6{,}704$&$8{,}964\,8063$&$7{,}962\,1126$&$6{,}721$&$8{,}964\,8449$&$7{,}962\,1288$\\
$6{,}671$&$559$&$7{,}961\,8106$&$6{,}688$&$721$&$7{,}962\,0073$&$6{,}705$&$119$&$165$&$6{,}722$&$434$&$263$\\
$6{,}672$&$702$&$244$&$6{,}689$&$827$&$161$&$6{,}706$&$170$&$200$&$6{,}723$&$413$&$235$\\
$6{,}673$&$844$&$381$&$6{,}690$&$8{,}964\,6930$&$245$&$6{,}707$&$219$&$232$&$6{,}724$&$386$&$201$\\
$6{,}674$&$8{,}964\,4984$&$514$&$6{,}691$&$8{,}964\,7031$&$326$&$6{,}708$&$263$&$260$&$6{,}725$&$354$&$164$\\
$6{,}675$&$8{,}964\,5121$&$642$&$6{,}692$&$130$&$404$&$6{,}709$&$303$&$284$&$6{,}726$&$316$&$123$\\
$6{,}676$&$257$&$769$&$6{,}693$&$225$&$478$&$6{,}710$&$338$&$305$&$6{,}727$&$273$&$079$\\
$6{,}677$&$391$&$7{,}961\,8892$&$6{,}694$&$316$&$552$&$6{,}711$&$370$&$321$&$6{,}728$&$226$&$7{,}962\,1030$\\
$6{,}678$&$522$&$7{,}961\,9012$&$6{,}695$&$405$&$622$&$6{,}712$&$399$&$336$&$6{,}729$&$173$&$7{,}962\,0977$\\
$6{,}679$&$652$&$131$&$6{,}696$&$491$&$690$&$6{,}713$&$423$&$345$&$6{,}730$&$114$&$920$\\
$6{,}680$&$780$&$246$&$6{,}697$&$573$&$755$&$6{,}714$&$441$&$351$&$6{,}731$&$8{,}964\,8051$&$\raisebox{-.3ex}{.}\;860$\\
$6{,}681$&$8{,}964\,5905$&$359$&$6{,}698$&$653$&$818$&$6{,}715$&$454$&$353$&$6{,}732$&$8{,}964\,7981$&$795$\\
$6{,}682$&$8{,}964\,6028$&$469$&$6{,}699$&$729$&$877$&$6{,}716$&$463$&$351$&$6{,}733$&$906$&$725$\\
$6{,}683$&$150$&$576$&$6{,}700$&$803$&$934$&$6{,}717$&$465$&$346$&$6{,}734$&$826$&$649$\\
$6{,}684$&$269$&$681$&$6{,}701$&$874$&$7{,}962\,0987$&$6{,}718$&$466$&$337$&$6{,}735$&$739$&$567$\\
$6{,}685$&$386$&$784$&$6{,}702$&$8{,}964\,7941$&$7{,}962\,1037$&$6{,}719$&$464$&$325$&$6{,}736$&$647$&$478$\\
$6{,}686$&$500$&$884$&$6{,}703$&$8{,}964\,8004$&$083$&$6{,}720$&$458$&$309$&$6{,}737$&$547$&$382$\\
$6{,}687$&$8{,}964\,6612$&$7{,}961\,9980$&$6{,}704$&$8{,}964\,8063$&$7{,}962\,1126$&$6{,}721$&$8{,}964\,8449$&$7{,}962\,1288$&$6{,}738$&$8{,}964\,7442$&$7{,}962\,0280$\\
\end{tabular}\par\endgroup
\end{center}
\vfill\hfill$18^*$\hspace*{2cm}\par
\clearpage
\noindent140\hfill{\footnotesize\textsc{Bemerkungen.\quad Konforme Übertragung des Sphäroids auf den Kegelmantel.}}\par\bigskip
\begin{center}\textsc{Bemerkungen.}\end{center}
\begingroup\small
Die vorstehende Aufzeichnung ist zusammen mit den beiden Tabellen einem Notizbuch entnommen. Sie dürfte aus dem Jahr 1823 oder 1824 stammen.

In den beiden Formeln für $6q$ wurden die Ausdrücke für die Glieder dritter Ordnung geändert. Bei \textsc{Gauss} lauten sie
\begingroup\footnotesize
\[
-\frac{(1+ee-3(1+4ee+e^4)\mu\mu+15(1+ee)(1+ee)\mu^4-14e^4\mu^6)(1-ee\mu\mu)(1-\mu\mu)}{3(1-ee)^3\mu^4}\rho^3\qquad .
\]
\par\endgroup
und
\begingroup\footnotesize
\[
-\frac{(1+ee-(1+6ee+7e^4)\mu\mu+(9ee+21e^4)\mu^4-14e^4\mu^6)(1-ee\mu\mu)(1-\mu\mu)}{3(1-ee)^3\mu^4}\delta^3.
\]
\par\endgroup
In der ersten Tabelle wurden die Überschriften der einzelnen Spalten hinzugefügt.

Für die geodätische Linie gilt
\[
\frac{\cos\varphi\sin A}{\sqrt{(1-ee\sin\varphi^2)}}=\text{const};\qquad(A=\text{Azimut})
\]
daher gilt wegen
\[
\frac rm=\frac{a\cos\varphi}{\mu\sqrt{(1-ee\sin\varphi^2)}},
\]
\[
\frac rm\sin A=\text{const.}
\]
Differenziert man diese Gleichung, so ergibt sich:
\[
\operatorname{cotang}A.\mathrm dA=\frac{\mathrm dm}{m}-\frac{\mathrm dr}{r}=-\frac{\sin\varphi}{\mu}.\frac{\mathrm dr}{r}
\]
oder, da $\operatorname{cotang}A=\dfrac{\mathrm dr}{r\,\mathrm d\Theta}$ und $A=\zeta-\Theta$ gilt:
\[
\mathrm d\zeta=\left(1-\frac{\sin\varphi}{\mu}\right)\mathrm d\Theta.
\]
Die Azimutreduktion $\zeta^0-Z$ erhält man nach Band IV, S.~278, aus der Gleichung
\[
\zeta^0-Z=(\tfrac13l^0+\tfrac16l')s\ldots,
\]
wobei $l^0$ und $l'$ hier die Werte von $l=-\dfrac{\mathrm dm}{m\,\mathrm dr}\sin\zeta=-\dfrac{6q}{r}\sin\zeta$ an den Endpunkten sind. Setzt man für $\dfrac{\sin\zeta}{r}$ einen mittleren Wert $=\dfrac{\Theta'-\Theta^0}{s}$ an, so ergibt sich:
\[
\zeta^0-Z=-(2q^0+q')(\Theta'-\Theta^0).
\]
\hfill\textsc{Krüger, Börsch.}\par
\par\endgroup
\vfill\begin{center}\rule{.24\textwidth}{.4pt}\end{center}
\clearpage
% Unnumbered title in the position of printed 141.
\null\vfill
\begin{center}
{\LARGE KONFORME ABBILDUNG\\[8pt]DES SPHÄROIDS IN DER EBENE.}\par\bigskip
{\Large (PROJEKTIONSMETHODE\\DER HANNOVERSCHEN LANDESVERMESSUNG.)}
\end{center}
\vfill\null
\endgroup
